% 【理论扩展】所属：第2章（由 chapters/revised/02-skepticism-luck.tex \input 引入）
\minitopic{怀疑论的范围：局部与全局}怀疑可以针对特定领域。道德怀疑论质疑道德知识，记忆怀疑论关注过去，自我知识怀疑论关注动机；这些都是\term{局部怀疑论}{local skepticism}。\term{全局怀疑论}{global skepticism}试图否定几乎所有经验知识。两者的论证负担不同：证明某类证据有系统缺陷，不能直接推出所有来源同样失败。 范围还涉及结论强度。“我们没有知识”“我们没有确定性”“我们的信念不合理”“我们无法向怀疑者证明”，是四个不同结论。无法给出不预设任何认识方法的证明，可能只显示彻底中立立场不可得，而不是普通知识不存在。阅读怀疑论时，应把否定词的对象和模态强度标出。 怀疑论也可能是方法而非结论。暂时搁置信念、主动寻找反例和压力测试来源，可以提高可靠性；永久要求排除逻辑上一切可能，则会使调查停滞。方法怀疑的合理限度取决于错误风险、核验成本与是否有具体反证。

\minitopic{古代怀疑与判断悬置}皮浪主义并不简单断言“知识不存在”，因为这本身像一个需要辩护的教条。怀疑者通过展示相反理由势均力敌，达到\term{判断悬置}{epoché}，并把宁静描述为随之而来的结果\parencite{sextus2000}。其挑战指向我们从表象跃迁到事物“本性”的资格。 阿格里帕式论证模式包括分歧、无穷回溯、相对性、假设和循环。任何主张似乎要么需要进一步理由而无穷后退，要么停在未经证明的假设，要么以自身或同一系统循环支持。这个五路难题比梦境更一般，因为它攻击理由结构，而非某一感官来源。 回应者可以接受可被击败的基础，认为理由不必由另一个信念支持；也可采用融贯结构，否认循环必然恶性；还可说分歧只在同侪与共享证据条件下产生强击败。每种回应都必须解释为何停止点不是任意的。

\minitopic{欠决定与最佳解释}若同一经验数据同时符合正常世界与完美模拟，数据似乎\term{欠决定}{underdetermine}两种理论。怀疑者由此说，我们没有理由偏向正常世界。回应者可能诉诸简单性、因果解释或先验概率：正常世界更直接解释经验的稳定互动，而模拟假设额外假定复杂装置。 怀疑者会反驳，我们凭什么在不知道外部结构时使用简单性？最佳解释若本身是待辩护认识规则，就可能把问题后移。回应并非毫无价值：它显示日常信念可能有相对解释优势；但它未必给出怀疑者要求的、完全不使用任何普通推理原则的证明。 欠决定在科学中通常是暂时和局部的。竞争理论会提出新预测，研究者寻找能区分它们的数据。全局怀疑情景被定义成与一切可能经验相容，因而原则上不提供区分预测。正因为它不可检验，其逻辑可能性未必应获得与普通解释相同权重。

\minitopic{铰链与理由的背景}维特根斯坦在《论确定性》中考察某些命题如何充当怀疑与求证得以进行的背景，例如世界已存在许久、物体不会在无人注意时任意消失\parencite{wittgenstein1969}。这些\term{铰链命题}{hinge propositions}未必是由更强证据推得的普通知识；它们像实践规则，使证据游戏可以运作。 铰链认识论可以说，要求同时证明所有背景承诺误解了理由的结构。我们在具体调查中固定大量背景，只有出现异常才局部修改。怀疑者则会追问：若铰链不由证据支持，把它们称为理性是否只是共同习惯？回应可强调它们并非任意内容，而是稳定行动、共同语言和纠错实践的条件。 铰链观点不应变成“凡是深信不疑就免于批评”。社会偏见也可能像背景一样不被注意，却应被改变。区别可能在于：真正的铰链是任何具体求证都预设的形式条件，偏见则可在保留调查实践时被证据揭露。具体界线仍有争议。

\minitopic{认识焦虑与两种要求}怀疑论常利用两种看似都合理的要求：知识必须向世界负责；主体必须能反思性地说明根据。外在主义通过可靠联系满足第一项，内在主义通过可及理由满足第二项。若坚持任何知识都必须同时以非循环方式满足二者，认识焦虑就会加剧\parencite{pritchard2016}。 一种双层方案区分动物知识与反思知识。正常知觉可因能力安全运行而产生一阶知识；主体理解能力何时可靠、能排除具体击败者时，获得更强的反思地位。反思层不需要证明所有认识来源总体可靠，否则回溯重现；它要求在相关范围内有负责任的视角。 这个方案解释了为什么儿童和动物可以知道，也解释了科学家为何有更高说明义务。代价是承认“知识”有层级或至少有不同认识成就。若日常语言只使用一个词，理论仍可区分其内部结构。

\minitopic{语境主义能否回答怀疑者}语境主义指出，普通语境中的“知道”可以为真，而怀疑讨论抬高标准后同一句话可能为假。它解释怀疑论何以在研讨室中有吸引力，又不迫使我们撤销全部日常归属。但怀疑者可能说，他问的是主体真实认识关系，而非词语在不同谈话中的伸缩。 因此语境主义更像语义诊断，而非直接证明外部世界。它能阻止从高标准语境结论无条件推广到普通语境，却仍需解释标准变化规则。若只要提到错误可能就能取消知识，怀疑者可以廉价操纵所有对话；若显著性受证据和会话目的约束，理论需要给出这些约束。 主体敏感的不变主义则让风险或证据影响主体是否知道，无需改变“知道”的语义。两种理论可能预测同一案例判断，却解释变化位置不同。语言调查、嵌入句测试和跨语境回指可作为区分证据。

\minitopic{局部怀疑的建设性作用}现实调查不需要先解决缸中脑，仍可提出局部怀疑。数字照片可能被编辑，所以要求元数据；药物研究可能有选择性发表，所以预注册；记忆可能受暗示，所以证人分开询问。这些制度把“可能出错”转化成具体纠错机制。 局部怀疑有三个优点。它说明错误机制，允许估计发生条件，提供可操作核验。全局怀疑只说一切证据可能整体欺骗，通常不给改进方向。认识实践可以承认无法消除全局逻辑可能，同时优先处理有机制、有迹象且后果重大的错误。

\minitopic{综合案例：远程身份验证}银行通过视频、证件照片和短信验证码确认客户身份。攻击者可能使用深度伪造、盗取手机并伪造证件。客户经理是否知道屏幕对面是本人？答案取决于知识标准和渠道独立性。三个步骤若都依赖同一被盗设备，不是三份独立证据；若视频活体检测、实体证件芯片与另一渠道回拨分别响应身份事实，组合更安全。 怀疑者可以设想攻击者完美控制一切渠道。银行不能逻辑排除，但可基于威胁模型、异常信号和损失限制达到可错知识与合理行动。高额转账提高复核要求，小额查询可使用较低门槛。这并不意味着身份事实随金额改变，而是证据是否足以承担行动的标准发生变化。 若事后发现是攻击者，不能只看错误结果倒推经理当时不理性；要检查当时可得信号和制度流程。若所有规范步骤都完成却遭遇前所未见攻击，信念可能合理但为假；若经理忽视明确警报，即使屏幕恰是真客户，也可能缺少信念确证。怀疑分析最终回到多维评价。

\begin{tcolorbox}[breakable,colback=white,colframe=blue!30!black,title={怀疑论诊断表},fonttitle=\bfseries]
写明怀疑对象；结论否定知识、确定性还是合理性；错误情景是逻辑可能还是有现实机制；论证是否使用闭合；标准是否要求不可错；回应拒绝哪个前提；该回应保存什么、牺牲什么；能否产生具体纠错行动。
\end{tcolorbox}
